\documentclass[10pt,a4paper]{amsart}
\usepackage[margin=19mm]{geometry}
\usepackage{amsmath,amssymb,mathtools}
\usepackage[scr=rsfs,cal=euler]{mathalfa}
\usepackage{xfrac}
\usepackage[unicode,hidelinks,bookmarks=false]{hyperref}
\newcommand{\ie}{i.e.\ }
\newcommand{\cf}{cf.\ }
\newcommand{\resp}{respectively\ }
\newcommand{\Subsection}{\subsection}
\providecommand{\nobreakdash}{\nobreak\hbox{-}}
\setlength{\emergencystretch}{2em}
\multlinegap0pt
\usepackage{amssymb}
\usepackage{mathtools}
\usepackage{dsfont}
\newtheorem{lemma}{Lemma}
\newcommand{\N}{\mathbb{N}}
\newcommand{\R}{\mathbb{R}}
\newcommand{\abs}[1]{\left|#1\right|}
\title{Bounds on Whittaker functions for $\mathrm{GL}(n)$}
\author{Gilles Felber, D\'avid T\'oth}
\date{Source: arXiv:2609.06501v1 (CC BY 4.0)}
\begin{document}
\maketitle

\noindent\emph{Source and licence.} Bounds on Whittaker functions for $\mathrm{GL}(n)$. Version arXiv:2609.06501v1 (CC BY 4.0), distributed by arXiv under the Creative Commons Attribution 4.0 International licence, http://creativecommons.org/licenses/by/4.0/, as recorded in the arXiv licence metadata for this exact version and shown on the abstract page https://arxiv.org/abs/2609.06501v1. Full licence notice: \texttt{licenses/arxiv-2609.06501-CC-BY-4.0.txt}.\par\smallskip

\noindent\textit{Editorial scope.} Counted unit: the article's lemma lemma (source0.tex lines 512--523). The article's complete original proof of it is delivered with it (source0.tex lines 524--550, one \texttt{proof} environment of 27 lines). No mathematics is written, repaired or re-derived: the statement and the proof are the article's own lines, in the article's own notation, and no retained line is substituted or re-flowed.  Retained uncounted as the source-local context the proof needs: lines 96--98; lines 245--254.  Nothing was omitted from any retained span, so the package carries no omission record.  No retained line contains a citation, so the wrapper carries no bibliography and no bibliography entry.  No retained line contains a figure environment, an image-inclusion command or drawing code, so no image is delivered and the package declares no asset dependency.  Editorial formatting only, in addition: an AMS \texttt{amsart} preamble that restates the article's own declarations and macros used by the retained lines, namely the environments \texttt{lemma} on separate counters, together with the standard packages those lines need.  The retained environments print in this document as Lemma 1 in order of appearance, whereas the article prints the same items with the numbering of its own full document.  The complete native source of the article as distributed in the arXiv e-print for this exact version is bundled unchanged as \texttt{sources/209541/source0.tex}.
\begin{equation}\label{eq:T_mu_def}
T_\mu:=\max(2,|\mu_1|,\dots,|\mu_n|).
\end{equation}

\begin{lemma}\label{lem bound integral of Bessel function}
Let $r\gg1$, $y>0$ and $\ell,k>0$ fixed.
Then we have
$$\int_0^\infty\abs{K_{ir}(2\pi y\sqrt{1+u^{\pm 2}})}^\ell u^{\pm k}\frac{du}u\ll_{\ell,k}e^{-\pi\ell r/2}r^{k-\ell/3}y^{-k}.$$
Moreover, if $2\pi y>1+\pi r/2$, then
$$\int_0^\infty\abs{K_{ir}(2\pi y\sqrt{1+u^{\pm 2}})}^\ell u^{\pm k}\frac{du}u\ll_{\ell,k}e^{-\pi\ell y}y^{-k-\ell/2}.$$
If in addition $r\gg T_\mu$ holds, where $T_\mu$ is defined in \eqref{eq:T_mu_def},
then we also have the uniform estimate
$$\int_0^\infty\abs{K_{ir}(2\pi y\sqrt{1+u^{\pm 2}})}^\ell u^{\pm k}\frac{du}u\ll_{\ell,k}e^{-\ell y/T_\mu}e^{-\pi\ell r/2}r^{k-\ell/3}y^{-k}.$$
\end{lemma}

\begin{lemma}
Let $n\in\N^+$, $r\geq1$ and $y_1,\dots,y_n>0$.
If
$\sum_{m=j}^nk_m$ is positive, then
\begin{align*}
\int_{\R_{>0}^n}K_{ir}&\left(2\pi y_1\sqrt{1+u_1^{-2}}\right)^2
\prod_{j=2}^nK_{ir}\left(2\pi y_j\sqrt{(1+u_{j-1}^2)(1+u_j^{-2})}\right)^2
\prod_{j=1}^nu_j^{-k_j}\frac{du_j}{u_j}\\
    &\ll_{k_1,\dots,k_n}
    \exp\left(-\pi nr-\frac2{T_\mu}\sum_{j=1}^n y_j\right)r^{\sum_{j=1}^n jk_j-2n/3}\prod_{j=1}^ny_j^{-\sum_{m=j}^nk_m}.
\end{align*}
\end{lemma}

\begin{proof}
We prove by induction on $n$. If $n=1$, then   Lemma~\ref{lem bound integral of Bessel function} 
gives the result.
Assume now that $n>1$ and the statements hold for $n-1$. We consider the integral over $u_n$. Again by Lemma \ref{lem bound integral of Bessel function}, we have
\[
\int_0^\infty K_{ir}\left(2\pi y_n\sqrt{(1+u_{n-1}^2)(1+u_n^{-2})}\right)^2u_n^{-k_n}\frac{du_n}{u_n}
\ll_{k_n}e^{-\pi r-\frac{2y_n}{T_\mu}}r^{k_n-\frac23}y_n^{-k_n}(1+u_{n-1}^2)^{-\frac{k_n}{2}}
\]
since by assumption $k_n>0$.

Now
we use $(1+u_{n-1}^2)^{-k_n/2}\leq u_{n-1}^{-k_n}$, and it remains
to estimate the integral 
\begin{align*}
I_{n-1}=\int_{\R_{>0}^{n-1}}&K_{ir}\left(2\pi y_1\sqrt{1+u_i^{-2}}\right)^2\prod_{j=2}^{n-1}K_{ir}\left(2\pi y_i\sqrt{(1+u_{j-1}^2)(1+u_j^{-2})}\right)^2\prod_{j=1}^{n-1}u_j^{-k_j'}\frac{du_j}{u_j},
\end{align*}
where $k_j'=k_j$ for $j=1,\dots,n-2$, while $k'_{n-1}=k_{n-1}+k_n$.  By
our original assumption we have 
$\sum_{m=j}^{n-1}k_m'>0$
for every $1\leq j\leq n-1$, so 
 the induction hypothesis gives
\begin{align*}
    I_{n-1}&\ll_{k_1,\dots,k_{n-1}}
    e^{-\pi(n-1)r-\frac2{T_\mu}\sum_{j=1}^{n-1}y_j}r^{\sum_{j=1}^{n-1} jk_j+(n-1)k_n-2(n-1)/3}\prod_{j=1}^{n-1}y_j^{-\sum_{m=j}^{n-1}k_m-k_n}
\end{align*}
and the statement follows.
\end{proof}

\end{document}
